% Canonical Paper I appendix.
This appendix supplements Proposition~\ref{prop:affine-cocycle-formulas} with
expanded calculations.  All discrete formulas hold over an arbitrary field
$K$, and every multiplier is assumed nonzero.

\subsection{Prefix recursion and closed form}

For the prefix history $(f_1,\ldots,f_k)$, write
\[
  F_k=f_k\circ\cdots\circ f_1,
  \qquad
  F_k(x)=\Phi_kx+\xi_k,
\]
with $F_0(x)=x$.  If $f_{k+1}(x)=s_{k+1}x+t_{k+1}$, then
\begin{equation}
  \Phi_{k+1}=s_{k+1}\Phi_k,
  \qquad
  \xi_{k+1}=s_{k+1}\xi_k+t_{k+1}.
  \label{eq:appendix-prefix-affine-recursion}
\end{equation}
Iterating the first recursion gives
\[
  \Phi_k=s_k\cdots s_1.
\]
Iterating the second one gives
\begin{align*}
  \xi_1&=t_1,\\
  \xi_2&=s_2t_1+t_2,\\
  \xi_3&=s_3s_2t_1+s_3t_2+t_3,\\
  \xi_4&=s_4s_3s_2t_1+s_4s_3t_2+s_4t_3+t_4.
\end{align*}
The pattern is
\[
  \xi_k=\sum_{i=1}^{k}t_i\prod_{j=i+1}^{k}s_j.
\]
Substitution in \eqref{eq:appendix-prefix-affine-recursion} proves this formula
inductively, including the final term through the empty-product convention.

The source-normalized prefix coordinate
$\widehat\xi_k=\Phi_k^{-1}\xi_k$ obeys a different recursion.  From
\eqref{eq:appendix-prefix-affine-recursion},
\begin{align}
  \widehat\xi_{k+1}
  &=\frac{s_{k+1}\xi_k+t_{k+1}}
          {s_{k+1}\Phi_k}
   =\widehat\xi_k+\frac{t_{k+1}}{\Phi_{k+1}}.
  \label{eq:appendix-normalized-prefix-recursion}
\end{align}
Starting from $\widehat\xi_0=0$ therefore gives
\[
  \widehat\xi_k
  =\sum_{i=1}^{k}\frac{t_i}{\Phi_i}
  =\sum_{i=1}^{k}\frac{t_i}{\prod_{j=1}^{i}s_j}.
\]
This second derivation makes the past normalization explicit at every prefix.

\subsection{Suffix transport and the future-weight kernel}

For $0\le i\le n$, define the scale of the strict suffix after step $i$ by
\[
  S_{i\to n}=\prod_{j=i+1}^{n}s_j,
  \qquad S_{n\to n}=1.
\]
Then the target-frame translation is simply
\begin{equation}
  \xi_n=\sum_{i=1}^{n}S_{i\to n}t_i.
  \label{eq:appendix-suffix-transport}
\end{equation}
To verify the interpretation, replace only $t_i$ by $t_i+\epsilon$.  All
states before step $i$ are unchanged.  At step $i$ the output gains
$\epsilon$, and every later affine map multiplies this change by its scale.
Thus the final endpoint changes by
\[
  \epsilon s_{i+1}\cdots s_n
  =\epsilon S_{i\to n}.
\]
The coefficient in \eqref{eq:appendix-suffix-transport} is therefore exactly
the transport of an additive event to the target frame.

\subsection{Alternating addition--scale histories}

Consider a positive real alternating history with pairs
\[
  x\xmapsto{\,e^{m_i}\,}e^{m_i}x
  \xmapsto{\,+\mu_i\,}e^{m_i}x+\mu_i,
  \qquad i=1,\ldots,n,
\]
applied from $i=1$ to $i=n$.  At the level of affine steps this has
$s_i=e^{m_i}$ and $t_i=\mu_i$.  Define prefix and suffix logarithmic
scales by
\[
  M_i=\sum_{j=1}^{i}m_j,
  \qquad
  M_{i+1:n}=\sum_{j=i+1}^{n}m_j.
\]
Proposition~\ref{prop:affine-cocycle-formulas} becomes
\begin{align}
  \Phi_n&=e^{M_n},
  \label{eq:appendix-alternating-scale}\\
  \xi_n&=\sum_{i=1}^{n}\mu_i e^{M_{i+1:n}},
  \label{eq:appendix-alternating-target}\\
  \widehat\xi_n&=\sum_{i=1}^{n}\mu_i e^{-M_i}.
  \label{eq:appendix-alternating-source}
\end{align}
For $n=3$, direct expansion checks the target formula:
\begin{align*}
  x
  &\longmapsto e^{m_1}x+\mu_1\\
  &\longmapsto e^{m_2+m_1}x
      +e^{m_2}\mu_1+\mu_2\\
  &\longmapsto e^{m_3+m_2+m_1}x
      +e^{m_3+m_2}\mu_1
      +e^{m_3}\mu_2+\mu_3.
\end{align*}
Dividing the translation line by
$e^{m_1+m_2+m_3}$ gives
\[
  \widehat\xi_3
  =\mu_1e^{-m_1}
   +\mu_2e^{-(m_1+m_2)}
   +\mu_3e^{-(m_1+m_2+m_3)},
\]
in agreement with \eqref{eq:appendix-alternating-source}.

The convention in this example is ``scale, then add'' within each pair.  If a
source instead groups the elementary operations as ``add, then scale'', the
translation attached to that pair is $e^{m_i}\mu_i$, not $\mu_i$.
Converting the pair to the common affine form $s_ix+t_i$ before using the
cocycle formulas prevents this convention change from being hidden.

\subsection{Concatenation in target and normalized coordinates}

Let
\[
  \rho(\gamma)=g(\Phi_\gamma,\xi_\gamma),
  \qquad
  \rho(\delta)=g(\Phi_\delta,\xi_\delta),
\]
where $g(\Phi,\xi)(x)=\Phi x+\xi$.  Since $\gamma\delta$ performs $\gamma$
first,
\begin{align*}
  \rho(\gamma\delta)(x)
  &=\rho(\delta)(\rho(\gamma)(x))\\
  &=\Phi_\delta(\Phi_\gamma x+\xi_\gamma)+\xi_\delta\\
  &=\Phi_\delta\Phi_\gamma x
    +\Phi_\delta\xi_\gamma+\xi_\delta.
\end{align*}
This yields the target-coordinate pair
\[
  (\Phi_{\gamma\delta},\xi_{\gamma\delta})
  =(\Phi_\delta\Phi_\gamma,
    \Phi_\delta\xi_\gamma+\xi_\delta).
\]
For normalized coordinates,
\begin{align*}
  \widehat\xi_{\gamma\delta}
  &=\frac{\Phi_\delta\xi_\gamma+\xi_\delta}
          {\Phi_\delta\Phi_\gamma}\\
  &=\widehat\xi_\gamma
    +\Phi_\gamma^{-1}\widehat\xi_\delta.
\end{align*}
The first formula transports the earlier translation forward by the later
scale.  The second transports the later normalized translation back through
the earlier scale.

For three histories, repeated use gives
\begin{align*}
  \xi_{\gamma\delta\epsilon}
  &=\Phi_\epsilon\Phi_\delta\xi_\gamma
    +\Phi_\epsilon\xi_\delta+\xi_\epsilon,\\
  \widehat\xi_{\gamma\delta\epsilon}
  &=\widehat\xi_\gamma
    +\Phi_\gamma^{-1}\widehat\xi_\delta
    +(\Phi_\delta\Phi_\gamma)^{-1}
      \widehat\xi_\epsilon.
\end{align*}
These are the block-history versions of future weighting and past
normalization.

\subsection{Inverse histories and relative comparison}

The affine inverse is
\begin{equation}
  g(\Phi,\xi)^{-1}
  =g(\Phi^{-1},-\Phi^{-1}\xi).
  \label{eq:appendix-affine-inverse}
\end{equation}
Indeed,
\begin{align*}
  g(\Phi^{-1},-\Phi^{-1}\xi)
  \bigl(\Phi x+\xi\bigr)
  &=x+\Phi^{-1}\xi-\Phi^{-1}\xi=x.
\end{align*}
Using \eqref{eq:appendix-affine-inverse}, the relative affine defect is
\begin{align}
  \rho(\delta)^{-1}\rho(\gamma)
  &=g\left(
      \frac{\Phi_\gamma}{\Phi_\delta},
      \frac{\xi_\gamma-\xi_\delta}{\Phi_\delta}
    \right).
  \label{eq:appendix-relative-affine-matrix}
\end{align}
In matrices, the same computation is
\[
  \begin{pmatrix}
    \Phi_\delta^{-1}&-\Phi_\delta^{-1}\xi_\delta\\0&1
  \end{pmatrix}
  \begin{pmatrix}
    \Phi_\gamma&\xi_\gamma\\0&1
  \end{pmatrix}
  =
  \begin{pmatrix}
    \Phi_\delta^{-1}\Phi_\gamma&
    \Phi_\delta^{-1}(\xi_\gamma-\xi_\delta)\\0&1
  \end{pmatrix}.
\]
If $\Phi_\gamma=\Phi_\delta=\Phi$, the relative element is a pure translation
by $(\xi_\gamma-\xi_\delta)/\Phi$.  The unnormalized endpoint difference is
$\xi_\gamma-\xi_\delta$.  These two scalar quantities differ by the common
scale and should not be identified.

\subsection{The elementary order square}

Let
\[
  A_\mu=
  \begin{pmatrix}1&\mu\\0&1\end{pmatrix},
  \qquad
  S_q=
  \begin{pmatrix}q&0\\0&1\end{pmatrix}.
\]
Under the chronological convention, first adding and then scaling has matrix
\[
  S_qA_\mu
  =\begin{pmatrix}q&q\mu\\0&1\end{pmatrix},
\]
whereas first scaling and then adding has matrix
\[
  A_\mu S_q
  =\begin{pmatrix}q&\mu\\0&1\end{pmatrix}.
\]
Their endpoint difference is
\[
  (qx+q\mu)-(qx+\mu)=\mu(q-1).
\]
Taking the second history as reference gives
\begin{align*}
  (A_\mu S_q)^{-1}(S_qA_\mu)
  &=\begin{pmatrix}1&\mu(1-q^{-1})\\0&1\end{pmatrix}.
\end{align*}
Thus the relative group translation is $\mu(1-q^{-1})$, while the open
endpoint defect is $\mu(q-1)$.  The latter is the former multiplied by the
common target scale $q$.

Closing the comparison into a group commutator produces yet another ordered
quantity.  For example,
\begin{align*}
  S_qA_\mu S_q^{-1}A_\mu^{-1}
  &=\begin{pmatrix}1&(q-1)\mu\\0&1\end{pmatrix},\\
  A_\mu S_qA_\mu^{-1}S_q^{-1}
  &=\begin{pmatrix}1&(1-q)\mu\\0&1\end{pmatrix}.
\end{align*}
The signs differ because the loop orientations differ.  These closed-loop
calculations are recorded only to audit conventions; the torsion theory and
its orientation choices are introduced in Section~\ref{sec:acs-torsion}.

\subsection{Exact perturbation propagation at every prefix}

Let $x$ and $\widetilde x$ be two initial values, and propagate both through
the same history.  Set
\[
  x_k=F_k(x),
  \qquad
  \widetilde x_k=F_k(\widetilde x).
\]
From the prefix recursion,
\begin{align*}
  \widetilde x_{k+1}-x_{k+1}
  &=\bigl(s_{k+1}\widetilde x_k+t_{k+1}\bigr)
    -\bigl(s_{k+1}x_k+t_{k+1}\bigr)\\
  &=s_{k+1}(\widetilde x_k-x_k).
\end{align*}
Induction gives
\begin{equation}
  \widetilde x_k-x_k
  =\left(\prod_{i=1}^{k}s_i\right)(\widetilde x-x)
  =\Phi_k(\widetilde x-x).
  \label{eq:appendix-prefix-perturbation}
\end{equation}
The translations cancel at each step; only the scale cocycle propagates the
initial difference.  Over $\mathbb R$, if $s_i=e^{\lambda_i}$, then
\[
  \frac{\widetilde x_k-x_k}{\widetilde x-x}
  =e^{\lambda_1+\cdots+\lambda_k}.
\]
No limiting process is involved in this identity.

\subsection{Differential conversion check}

Finally, on the positive real affine group let $\Phi=e^M$ and
$\widehat\xi=e^{-M}\xi$.  The product rule gives
\begin{align*}
  d\widehat\xi
  &=d(e^{-M}\xi)\\
  &=-e^{-M}\xi\,dM+e^{-M}d\xi\\
  &=e^{-M}(d\xi-\xi dM).
\end{align*}
Therefore
\[
  d\xi-\xi dM=e^M d\widehat\xi.
\]
This verifies that the translation component of $dg\,g^{-1}$ is the target
scale times the differential of the source-normalized coordinate.  It does
not identify $dg\,g^{-1}$ with $g^{-1}dg$, whose translation component is
$e^{-M}d\xi$.
